\documentclass[a4paper,12pt]{article}
\usepackage[T2A]{fontenc}
\usepackage[utf8]{inputenc}
\usepackage[russian]{babel}
\usepackage{amsmath,amssymb,mathtools}
\usepackage{helvet}
% Sans-serif T2A: outline CM-LGC glyphs for Cyrillic pdfLaTeX.
\renewcommand{\sfdefault}{fcs}
\renewcommand{\familydefault}{\sfdefault}
\usepackage{icomma}
\usepackage[a4paper,left=23mm,right=22mm,top=22mm,bottom=25mm]{geometry}
\usepackage{microtype}
\usepackage{tikz}
\usetikzlibrary{calc,arrows.meta,positioning,shapes.geometric}
\usepackage{pgfplots}
\pgfplotsset{compat=1.18}
\usepackage{booktabs,tabularx}
\usepackage{enumitem}
\usepackage{hyperref}
\usepackage{fancyhdr}
\usepackage{setspace}
\usepackage{xcolor}
\definecolor{accent}{HTML}{008080}
\definecolor{darkaccent}{HTML}{004D4D}
\definecolor{textgray}{HTML}{333333}
\definecolor{softgray}{HTML}{647277}
\hypersetup{colorlinks=true,linkcolor=darkaccent,urlcolor=darkaccent,hypertexnames=false}
\setstretch{1.25}
\setlength{\parindent}{0pt}
\setlength{\parskip}{6pt}
\setlist[itemize]{leftmargin=1.9em,itemsep=4pt,topsep=4pt}
\setlist[enumerate]{leftmargin=1.8em,itemsep=7pt,topsep=5pt}
\pagestyle{fancy}
\fancyhf{}
\renewcommand{\headrulewidth}{0pt}
\fancyfoot[C]{\footnotesize\color{softgray} Лёвин Артём Александрович эксклюзивно для Ярослава Гаврилова}
\fancyfoot[R]{\footnotesize\thepage}
\newcommand{\sectionline}{\par\vspace{2pt}\noindent\textcolor{accent}{\rule{\linewidth}{0.5pt}}\par\vspace{3pt}}
\newcommand{\focus}[1]{\textbf{\textcolor{darkaccent}{#1}}}
\newcommand{\checklist}[1]{\par\smallskip\textcolor{darkaccent}{\textbf{Проверьте себя.}}\par\begin{itemize}[label=$\square$]#1\end{itemize}}
\newcommand{\tgfun}{\operatorname{tg}}
\newcommand{\ctgfun}{\operatorname{ctg}}
\newcommand{\arctgfun}{\operatorname{arctg}}
\begin{document}
\begin{titlepage}
\thispagestyle{empty}
\begin{center}
\vspace*{40mm}
{\large\color{softgray} МАТЕМАТИКА \enspace / \enspace ПРОФИЛЬНЫЙ ЕГЭ}\par
\vspace{13mm}
{\fontsize{36}{42}\selectfont\bfseries\color{darkaccent} Чек-лист}\par
\vspace{7mm}
{\LARGE\bfseries Степени, прикладные задачи\par и показательная функция}\par
\vspace{14mm}
{\large Ярослав Гаврилов}\par
\vspace{3mm}
{\large 22 октября 2026 года}\par
\vspace{15mm}
{\large\color{textgray} Правила \(\longrightarrow\) разобранные примеры \\ \(\longrightarrow\) самостоятельная тренировка}\par
\vfill
{\normalsize\color{softgray} Лёвин Артём Александрович}\par
{\normalsize\color{softgray} эксклюзивно для Ярослава Гаврилова}\par
\end{center}
\begin{tikzpicture}[remember picture,overlay]
\draw[accent!58,line width=1.2pt] 
([xshift=23mm,yshift=12mm]current page.south west)
.. controls ([xshift=69mm,yshift=25mm]current page.south west) and
([xshift=-70mm,yshift=3mm]current page.south east) ..
([xshift=-22mm,yshift=15mm]current page.south east);
\end{tikzpicture}
\end{titlepage}
\setcounter{page}{1}
\section*{1. Степени: инструменты для вычислений}
\sectionline
Для решения задач со степенями удобно \focus{сначала представить числа через степени}, затем выполнить преобразования и только в конце переходить к вычислениям.
\[
\begin{aligned}
 a^m\cdot a^n&=a^{m+n}, &\frac{a^m}{a^n}&=a^{m-n},\\
 (a^m)^n&=a^{mn}, &(ab)^n&=a^n b^n,\\
 a^{-n}&=\frac{1}{a^n}, &a^{1/n}&=\sqrt[n]{a}.
\end{aligned}
\]
Здесь в формулах с дробями основания и знаменатели не равны нулю. Для дробного показателя \(1/n\) и чётного \(n\) требуется \(a\ge 0\); при работе с \(a^x\) для произвольного действительного \(x\) берут \(a>0\).

\textbf{Пример.} Упростите выражение:
\[
\frac{10^{21}}{10^{-9}\cdot10^{20}}
=10^{21-(-9)-20}=10^{10}.
\]
При делении показатель степени \focus{вычитают целиком}: \(21-(-9)-20=10\).

\textbf{Пример.} Разложите коэффициент на простые множители:
\[
15625=5^6,\qquad 228=2^2\cdot3\cdot19,\qquad 57=3\cdot19,
\qquad 10^{10}=2^{10}\cdot5^{10}.
\]

\checklist{
\item Умеете приводить произведения и частные степеней к одному основанию.
\item Различаете \(10^{-9}\) и \(-10^9\); помните про двойной минус в показателе.
\item Умеете переходить от десятичной записи к целому коэффициенту: \(5{,}7\cdot10^{-8}=57\cdot10^{-9}\).
\item Перед извлечением корня проверяете знак числа и смысл искомой величины.
}

\section*{2. Показательные уравнения и распад}
\sectionline
При \(a>0\), \(a\ne 1\) из \(a^u=a^v\) следует \(u=v\).

\textbf{Пример.} Масса вещества изменяется по закону
\[
m(t)=m_0\cdot 2^{-t/T},\qquad m_0=40,\quad T=10.
\]
Найдите время \(t\), когда масса станет равной \(5\).
\[
5=40\cdot2^{-t/10}\quad\Longleftrightarrow\quad
2^{-t/10}=\frac18=2^{-3}\quad\Longleftrightarrow\quad
t=30.
\]
Запись \(m(t)\) означает \focus{значение массы в момент \(t\)}; это запись функции, а не произведение \(m\cdot t\).

\checklist{
\item Умеете из формулы выделять степень с неизвестным показателем.
\item Представляете дроби вида \(1/8\), \(1/16\), \(1/27\) как степени.
\item Проверяете знак ответа: время распада не может быть отрицательным.
}

\section*{3. Два состояния: давление и объём газа}
\sectionline
Если \(pV^a=C\), то для двух состояний одного газа выполняется
\[
p_1V_1^a=p_2V_2^a,\qquad p_1,p_2,V_1,V_2>0.
\]
\textbf{Пример.} Объём уменьшился вдвое, давление увеличилось вчетверо. Найдите \(a\).
\[
V_2=\frac{V_1}{2}\quad\Longleftrightarrow\quad V_1=2V_2,
\qquad p_2=4p_1.
\]
\[
\begin{aligned}
p_1(2V_2)^a&=(4p_1)V_2^a,\\
p_1\cdot2^a\cdot V_2^a&=4p_1V_2^a,\\
2^a&=4=2^2,\qquad \boxed{a=2}.
\end{aligned}
\]
Сокращение допустимо, поскольку \(p_1>0\) и \(V_2^a>0\). При подстановке \focus{составного выражения} обязательно ставьте скобки: \((2V_2)^a\).

\checklist{
\item Умеете переводить «уменьшился в \(k\) раз» в \(V_2=V_1/k\).
\item Составляете равенство для состояний «до» и «после».
\item Корректно возводите произведение в степень и сокращаете общий ненулевой множитель.
}

\section*{4. Степень неизвестной: извлечение корня}
\sectionline
\textbf{Пример.} Сила натяжения троса определяется формулой
\[
F=\rho gL^3-mg.
\]
Пусть \(\rho=1000\), \(g=10\), \(m=25\). Найдите наибольшее физически допустимое \(L\ge0\), при котором \(F\le1000\).
\[
10000L^3-250\le1000
\quad\Longleftrightarrow\quad L^3\le\frac{1250}{10000}=\frac18.
\]
Так как функция \(L^3\) возрастает, то
\[
L\le\sqrt[3]{\frac18}=\frac12.
\qquad\boxed{L_{\max}=0{,}5}.
\]
\focus{Пограничное значение} находят из равенства \(F=1000\), а знак неравенства сохраняют, чтобы обосновать, почему найденная величина максимальна.

\textbf{Правило.} Если \(x^n=A\) и \(n\) нечётно, то \(x=\sqrt[n]{A}\). Если \(n\) чётно и \(A>0\), уравнение имеет два действительных корня \(x=\pm\sqrt[n]{A}\). Для длины и абсолютной температуры отрицательное значение исключается по смыслу.

\checklist{
\item Различаете \(L\) и \(L^3\): после нахождения \(L^3\) извлекаете кубический корень.
\item Понимаете, почему «не более» означает \(\le\), а «не менее» --- \(\ge\).
\item Учитываете ограничения физической задачи при выборе корня.
}

\clearpage
\section*{Алгоритм: как решать прикладную задачу со степенями}
\vspace{3mm}
\begin{center}
\begin{tikzpicture}[
 node distance=9mm and 13mm,
 >=Stealth,
 start/.style={ellipse,draw=darkaccent,thick,align=center,minimum width=46mm,minimum height=10mm,font=\small},
 action/.style={rectangle,rounded corners=2pt,draw=accent,thick,align=center,text width=102mm,inner sep=6pt,font=\small},
 test/.style={diamond,aspect=3.2,draw=darkaccent,thick,align=center,text width=40mm,inner sep=1.5pt,font=\small},
 arrow/.style={->,>=Stealth,line width=0.9pt,draw=darkaccent}
]
\node[start] (start) {Начало};
\node[action,below=of start] (data) {Выпишите формулу, данные, искомую величину и ограничения};
\node[test,below=of data] (states) {Есть два состояния?};
\node[action,below left=13mm and 8mm of states,text width=43mm] (two) {Запишите связь\\[2pt]$p_1V_1^a=p_2V_2^a$};
\node[action,below right=13mm and 8mm of states,text width=43mm] (one) {Подставьте данные в одну формулу};
\node[action,below=32mm of states] (isolate) {Упростите коэффициенты и оставьте искомую степень отдельно};
\node[test,below=of isolate] (unknown) {Неизвестная в показателе?};
\node[action,below left=13mm and 8mm of unknown,text width=43mm] (exp) {Приведите к одному основанию и приравняйте показатели};
\node[action,below right=13mm and 8mm of unknown,text width=43mm] (root) {Выделите \(x^n\) и извлеките допустимый корень};
\node[action,below=35mm of unknown] (check) {Проверьте вычисления, ограничения и смысл результата};
\node[start,below=of check] (end) {Ответ};
\draw[arrow] (start)--(data);
\draw[arrow] (data)--(states);
\draw[arrow] (states.south west) -| node[pos=0.26,above,sloped,font=\scriptsize]{да} (two.north);
\draw[arrow] (states.south east) -| node[pos=0.26,above,sloped,font=\scriptsize]{нет} (one.north);
\draw[arrow] (two.south) |- ([yshift=8mm]isolate.north) -- (isolate.north);
\draw[arrow] (one.south) |- ([yshift=8mm]isolate.north);
\draw[arrow] (isolate)--(unknown);
\draw[arrow] (unknown.south west) -| node[pos=0.25,above,sloped,font=\scriptsize]{да} (exp.north);
\draw[arrow] (unknown.south east) -| node[pos=0.25,above,sloped,font=\scriptsize]{нет} (root.north);
\draw[arrow] (exp.south) |- ([yshift=7mm]check.north) -- (check.north);
\draw[arrow] (root.south) |- ([yshift=7mm]check.north);
\draw[arrow] (check)--(end);
\end{tikzpicture}
\end{center}

\clearpage
\section*{5. Большие числа и четвёртая степень: температура звезды}
\sectionline
Закон Стефана--Больцмана:
\[
P=\sigma S T^4,\qquad\sigma>0,\quad S>0,\quad T>0.
\]
\textbf{Пример.} Даны
\[
\sigma=5{,}7\cdot10^{-8},\qquad
S=\frac{10^{20}}{228},\qquad
P\ge 1{,}5625\cdot10^{25}.
\]
Определите наименьшую возможную температуру \(T\) (в кельвинах).

\textbf{Шаг 1.} Для минимальной температуры используем граничное равенство и выделяем \(T^4\):
\[
T^4=\frac{P}{\sigma S}.
\]
Здесь \(P=\sigma ST^4\) возрастает по \(T>0\), поэтому равенство действительно даёт минимум.

\textbf{Шаг 2.} Убираем десятичные дроби и оформляем одну дробь:
\[
5{,}7\cdot10^{-8}=57\cdot10^{-9},\qquad
1{,}5625\cdot10^{25}=15625\cdot10^{21},
\]
\[
T^4=\frac{15625\cdot10^{21}\cdot228}
{57\cdot10^{-9}\cdot10^{20}}.
\]
\textbf{Шаг 3.} Раскладываем и сокращаем множители:
\[
\begin{aligned}
T^4
&=\frac{5^6\cdot10^{21}\cdot(2^2\cdot3\cdot19)}
{(3\cdot19)\cdot10^{-9}\cdot10^{20}}\\
&=5^6\cdot2^2\cdot10^{21-(-9)-20}\\
&=5^6\cdot2^2\cdot(2\cdot5)^{10}
=2^{12}\cdot5^{16}.
\end{aligned}
\]
\textbf{Шаг 4.} Извлекаем корень четвёртой степени:
\[
T=(2^{12}\cdot5^{16})^{1/4}
=2^3\cdot5^4=8\cdot625=\boxed{5000\ \text{К}}.
\]
\focus{Проверка.} \(5000^4=6{,}25\cdot10^{14}\); подстановка в \(P=\sigma ST^4\) возвращает \(1{,}5625\cdot10^{25}\) Вт.

\checklist{
\item Переносите десятичную запятую согласованно с показателем десяти.
\item Понимаете, почему \(\frac{10^{21}}{10^{-9}\cdot10^{20}}=10^{10}\).
\item Можете восстановить все шаги факторизации и извлечения корня.
\item Проверяете последний числовой результат, даже если преобразования степеней правильны.
}

\section*{6. Показательная функция по графику}
\sectionline
Рассмотрим \(f(x)=a^x+b\), где \(a>0\), \(a\ne1\). По графику видны точки \(A(0;-2)\) и \(B(1;-1)\).
\begin{center}
\begin{tikzpicture}
\begin{axis}[
width=0.78\linewidth,height=62mm,
axis lines=middle, axis line style={->,black},
xmin=-1.6,xmax=3.2,ymin=-3.5,ymax=6.2,
xtick={-1,0,1,2,3},ytick={-3,-2,-1,0,1,3,5},
xlabel={$x$},ylabel={$y$},
clip=false, grid=major, grid style={gray!18},
tick label style={font=\small},xlabel style={at={(axis description cs:0.99,0.23)},anchor=east},
ylabel style={at={(axis description cs:0.48,0.99)},anchor=north}]
\addplot[accent,thick,samples=220,domain=-1.5:3.1]{2^x-3};
\addplot[only marks,mark=*,mark size=2.4pt,black] coordinates {(0,-2) (1,-1)};
\node[anchor=north west,font=\small] at (axis cs:0.10,-2.08) {$A$};
\node[anchor=north west,font=\small] at (axis cs:1,-1) {$B$};
\end{axis}
\end{tikzpicture}
\end{center}
\[
\begin{cases}
-2=a^0+b=1+b,\\
-1=a^1+b=a+b
\end{cases}
\quad\Longrightarrow\quad
b=-3,\ a=2.
\]
Следовательно,
\[
f(x)=2^x-3,\qquad f(6)=64-3=\boxed{61}.
\]
\checklist{
\item Узнаёте, что \(a^0=1\), и находите сдвиг \(b\) по точке с \(x=0\).
\item Определяете \(a\) по второй точке и затем вычисляете \(f(x)\).
\item Учитываете: при \(a>1\) график возрастает, при \(0<a<1\) --- убывает.
}

\clearpage
\section*{Тренировочный блок \hfill Путь А}
\sectionline
\textcolor{softgray}{10 задач. Решайте по порядку; в задачах 5--10 записывайте основные преобразования, а не только ответ.}
\begin{enumerate}
\item Решите уравнение \(2^{3x-5}=\dfrac1{16}\).
\item Найдите \(t\ge0\), если \(6=48\cdot2^{-t/12}\).
\item Вычислите \(\dfrac{10^{14}}{10^{-4}\cdot10^{12}}\).
\item Вычислите \(\dfrac{3^4\cdot9^2}{27}\).
\item Давление газа и его объём связаны законом \(pV^a=C\). При уменьшении объёма втрое давление выросло в 27 раз. Найдите \(a\).
\item Время полураспада равно 8 дням, начальная масса равна 96 г. Через сколько дней масса, изменяющаяся по формуле \(m(t)=96\cdot2^{-t/8}\), станет равной 12 г?
\item Сила определяется формулой \(F=\rho gL^3-mg\), где \(\rho=1000\), \(g=10\), \(m=20\). Найдите наибольшее \(L\ge0\), при котором \(F\le1050\).
\item График функции \(f(x)=a^x+b\) проходит через точки \((0;-4)\) и \((1;-2)\). Найдите \(f(4)\) (\(a>0\), \(a\ne1\)).
\item Мощность излучения равна \(P=\sigma ST^4\). Известно, что \(\sigma=5{,}7\cdot10^{-8}\), \(S=\dfrac{10^{20}}{64}\), \(P\ge2{,}28\cdot10^{25}\). Найдите наименьшую температуру \(T>0\) в кельвинах.
\item Найдите положительное число \(x\), если
\[
x^4=\frac{15625\cdot228\cdot10^{17}}{57\cdot10^{-3}\cdot10^{14}}.
\]
Представьте все множители в виде степеней и приведите вычисление к простым основаниям.
\end{enumerate}

\clearpage
\section*{Ответы к тренировочному блоку}
\sectionline
\renewcommand{\arraystretch}{1.55}
\begin{center}
\begin{tabularx}{0.88\linewidth}{@{}cX@{}}
\toprule
\textbf{№} & \textbf{Ответ} \\
\midrule
1 & \(x=\dfrac13\) \\
2 & \(t=36\) \\
3 & \(10^6=1\,000\,000\) \\
4 & \(243\) \\
5 & \(a=3\) \\
6 & \(24\) дня \\
7 & \(L_{\max}=0{,}5\) \\
8 & \(f(4)=76\) \\
9 & \(4000\) К \\
10 & \(x=500\) \\
\bottomrule
\end{tabularx}
\end{center}
\vspace{9mm}
\textbf{Самоконтроль по сложным заданиям.}

В задаче 7 должно получиться \(L^3\le\frac18\). В задаче 9 при граничной мощности получится \(T^4=2{,}56\cdot10^{14}=4000^4\). В задаче 10 перед извлечением корня получится \(x^4=2^8\cdot5^{12}\).

\vfill
\begin{center}
\textcolor{darkaccent}{\rule{.5\linewidth}{.5pt}}\\[6pt]
\small Повторите задания, в которых понадобилась подсказка или был потерян знак степени.
\end{center}
\end{document}
